\chapter{جزئیات کالیبراسیون دوربین استریو}
\label{app:stereo_calibration_details}

\section{هدف از کالیبراسیون دوربین استریو}
\label{app:stereo_calibration_objective}

در سامانه عملی این پژوهش، تخمین فاصله هدف و موانع بر پایه پردازش تصاویر دوربین استریو انجام می‌شود. از آنجا که محاسبه عمق در سامانه بینایی استریو به پارامترهای هندسی دوربین وابسته است، پیش از استفاده از تصاویر در زنجیره ادراک، فرایند کالیبراسیون دوربین استریو انجام شد. هدف از این کالیبراسیون، محاسبه پارامترهای درونی هر دوربین، ضرایب اعوجاج لنز، رابطه هندسی میان دو دوربین و ماتریس‌های لازم برای تصحیح هندسی تصاویر چپ و راست است.

در حالت ایده‌آل، پس از تصحیح هندسی، نقاط متناظر در تصاویر چپ و راست باید روی خط‌های افقی متناظر قرار گیرند. این ویژگی باعث می‌شود جست‌وجوی نقاط متناظر در الگوریتم‌های تطبیق استریو تنها در راستای افقی انجام شود و محاسبه اختلاف مکانی پایداری بیشتری داشته باشد. در صورتی که کالیبراسیون یا ترتیب دوربین‌های چپ و راست نادرست باشد، نقشه اختلاف مکانی دچار پراکندگی، ناپیوستگی و خطاهای شدید خواهد شد. بنابراین، کیفیت کالیبراسیون مستقیماً بر دقت فاصله‌سنجی استریو اثر می‌گذارد.

\section{دریافت تصاویر مورد استفاده برای کالیبراسیون}
\label{app:stereo_calibration_image_capture}

برای ثبت تصاویر کالیبراسیون، دوربین استریو \lr{IMX219-83} در همان وضعیت نصب‌شده روی سامانه عملی مورد استفاده قرار گرفت. این موضوع اهمیت دارد، زیرا هرگونه تغییر در موقعیت نسبی دو دوربین، زاویه نصب، جهت تصویر یا فاصله میان دوربین‌ها پس از کالیبراسیون می‌تواند اعتبار پارامترهای محاسبه‌شده را از بین ببرد. بنابراین، کالیبراسیون با همان چینش فیزیکی و همان جهت تصویری انجام شد که در سامانه عملی برای دریافت تصویر استفاده می‌شود.

تصاویر خام توسط رزبری پای ۵ دریافت و از طریق ارتباط شبکه به لپ‌تاپ منتقل شدند. در سمت رزبری پای، تصاویر با اندازه $1640 \times 1232$ و در قالب \lr{YUV420} دریافت شدند. از این قالب، تنها صفحه روشنایی یا \lr{Y-plane} استخراج شد؛ زیرا این صفحه متناظر با تصویر خاکستری کامل است و برای کالیبراسیون و پردازش استریو اطلاعات شدت روشنایی کافی را فراهم می‌کند. همچنین برای جلوگیری از ناسازگاری با نقشه‌های تصحیح هندسی، تصویر استخراج‌شده دقیقاً به اندازه کالیبراسیون برش داده شد.

در کد دریافت تصویر، جهت تصویر به‌صورت ثابت نگه داشته شد و از اعمال چرخش یا وارون‌سازی نرم‌افزاری پس از کالیبراسیون خودداری گردید. اگرچه در برخی تصاویر، جهت نمایش ظاهری تصویر وارونه به نظر می‌رسد، این وضعیت از نظر هندسی مشکلی ایجاد نمی‌کند؛ زیرا تمام مراحل کالیبراسیون، تصحیح هندسی و تولید اختلاف مکانی با همان جهت ثابت انجام شده‌اند. تغییر جهت تصویر پس از کالیبراسیون باعث برهم‌خوردن هم‌راستاسازی هندسی تصاویر و خطا در محاسبه اختلاف مکانی می‌شود.

برای بهبود کیفیت تصاویر کالیبراسیون، نوردهی دو دوربین تا حد امکان یکسان تنظیم شد. در تصویربرداری استریو، تفاوت شدید نوردهی یا بهره بین دوربین چپ و راست می‌تواند باعث شود یک نقطه فیزیکی در دو تصویر شدت روشنایی متفاوتی داشته باشد و در نتیجه، آشکارسازی گوشه‌های شطرنجی یا تطبیق استریو با خطا همراه شود. به همین دلیل، در زمان ثبت تصاویر، تنظیمات نوردهی و بهره دوربین‌ها کنترل شد تا تصاویر چپ و راست از نظر روشنایی و کنتراست تا حد امکان سازگار باشند.

\section{الگوی کالیبراسیون و نحوه ثبت نمونه‌ها}
\label{app:stereo_calibration_board_capture}

برای کالیبراسیون، از الگوی شطرنجی با ابعاد مشخص استفاده شد. مشخصات الگو به‌صورت \lr{10x7x25} تعریف شد؛ به این معنا که تعداد گوشه‌های داخلی الگو برابر $10 \times 7$ و اندازه هر خانه شطرنجی برابر $25~mm$ در نظر گرفته شد. بر اساس این مشخصات، مختصات نقاط متناظر در دستگاه مختصات واقعی ساخته شد و سپس با مختصات دوبعدی گوشه‌های آشکارسازی‌شده در تصاویر چپ و راست متناظر گردید.

در هنگام ثبت تصاویر، الگوی شطرنجی در موقعیت‌ها، فاصله‌ها و زاویه‌های مختلف نسبت به دوربین قرار داده شد. ثبت تصویر تنها از مرکز میدان دید یا تنها در یک فاصله ثابت برای کالیبراسیون کافی نیست؛ زیرا در این حالت، پارامترهای دوربین عمدتاً بر اساس ناحیه‌ای محدود از تصویر تخمین زده می‌شوند. به همین دلیل، در مجموعه تصاویر کالیبراسیون، تلاش شد الگو در بخش‌های مختلف میدان دید، در فاصله‌های نزدیک و دور، و با میزان‌های مختلف چرخش و کجی قرار گیرد.

در نرم‌افزار کالیبراسیون، برای کنترل تنوع داده‌های ثبت‌شده، چند معیار پوشش در نظر گرفته شد. این معیارها شامل پوشش افقی الگو در تصویر، پوشش عمودی، تغییر اندازه ظاهری الگو و میزان چرخش یا کجی الگو بودند. این شاخص‌ها کمک می‌کنند مجموعه تصاویر تنها شامل حالت‌های تکراری نباشد و برای تخمین پارامترهای کالیبراسیون، داده‌هایی با تنوع هندسی کافی در اختیار الگوریتم قرار گیرد.

نمونه‌هایی از تصاویر ثبت‌شده برای کالیبراسیون در شکل \ref{fig:stereo_calibration_image_samples} نشان داده شده است. همان‌طور که مشاهده می‌شود، الگوی شطرنجی در حالت‌های مختلف نسبت به دوربین قرار گرفته است. تصاویری که در آن‌ها الگو در هر دو تصویر چپ و راست به‌طور کامل قابل تشخیص نبود، یا تصویر دارای تاری محسوس بود، برای کالیبراسیون استفاده نشدند.

\begin{figure}[htbp]
	\centering
	
	\begin{minipage}{0.48\linewidth}
		\centering
		\includegraphics[width=\linewidth]{left_001.png}
		\vspace{2mm}
		
		{\small نمونه تصویر چپ، حالت اول}
	\end{minipage}
	\hfill
	\begin{minipage}{0.48\linewidth}
		\centering
		\includegraphics[width=\linewidth]{right_001.png}
		\vspace{2mm}
		
		{\small نمونه تصویر راست، حالت اول}
	\end{minipage}
	
	\vspace{4mm}
	
	\begin{minipage}{0.48\linewidth}
		\centering
		\includegraphics[width=\linewidth]{left_035.png}
		\vspace{2mm}
		
		{\small نمونه تصویر چپ، حالت دوم}
	\end{minipage}
	\hfill
	\begin{minipage}{0.48\linewidth}
		\centering
		\includegraphics[width=\linewidth]{right_035.png}
		\vspace{2mm}
		
		{\small نمونه تصویر راست، حالت دوم}
	\end{minipage}
	
	\caption{نمونه‌هایی از تصاویر جفتی ثبت‌شده برای کالیبراسیون دوربین استریو}
	\label{fig:stereo_calibration_image_samples}
\end{figure}

\section{آشکارسازی گوشه‌ها و محاسبه پارامترهای کالیبراسیون}
\label{app:stereo_calibration_computation}

پس از ثبت جفت‌تصویرهای معتبر، گوشه‌های داخلی الگوی شطرنجی در هر دو تصویر چپ و راست آشکارسازی شدند. برای هر نمونه معتبر، مختصات دوبعدی گوشه‌ها در تصویر چپ و راست با مختصات سه‌بعدی متناظر روی صفحه الگوی شطرنجی مرتبط شد. این داده‌ها ورودی اصلی فرایند کالیبراسیون را تشکیل می‌دهند.

فرایند کالیبراسیون در دو مرحله انجام شد. در مرحله نخست، هر دوربین به‌صورت مستقل کالیبره شد. در این مرحله، ماتریس درونی دوربین چپ، ماتریس درونی دوربین راست و ضرایب اعوجاج هر دوربین محاسبه شدند. ماتریس درونی شامل فاصله‌های کانونی مؤثر و مختصات مرکز تصویر است. ضرایب اعوجاج نیز برای مدل‌سازی و اصلاح اعوجاج‌های ناشی از لنز استفاده می‌شوند.

در مرحله دوم، کالیبراسیون استریو با استفاده از نقاط متناظر میان دو تصویر انجام شد. در این مرحله، رابطه هندسی میان دو دوربین شامل ماتریس دوران و بردار انتقال محاسبه گردید. مؤلفه اصلی بردار انتقال در راستای افقی، فاصله پایه میان دو دوربین را مشخص می‌کند. این فاصله پایه همراه با فاصله کانونی تصویر تصحیح‌شده در رابطه تبدیل اختلاف مکانی به عمق استفاده می‌شود؛ بنابراین، خطا در تخمین این پارامترها مستقیماً بر دقت فاصله‌سنجی اثر می‌گذارد.

پس از محاسبه پارامترهای درونی و رابطه نسبی دو دوربین، ماتریس‌های لازم برای تصحیح هندسی تصاویر محاسبه شدند. در این مرحله، تصاویر چپ و راست به‌گونه‌ای بازنگاشت می‌شوند که اثر اعوجاج کاهش یافته و خطوط متناظر دو تصویر هم‌راستا شوند. خروجی این مرحله شامل ماتریس‌های تصحیح، ماتریس‌های تصویر پس از تصحیح، ماتریس \lr{Q} و ناحیه معتبر تصویر است.

خروجی نهایی کالیبراسیون در قالب فایل \lr{JSON} ذخیره شد. این فایل شامل ماتریس‌های درونی دوربین‌ها، ضرایب اعوجاج، ماتریس دوران، بردار انتقال، ماتریس‌های تصحیح، ماتریس‌های تصویر پس از تصحیح، ماتریس \lr{Q}، ناحیه معتبر تصویر، اندازه تصویر، خطای بازفرافکنی و فاصله پایه دوربین‌ها است. این فایل در ادامه در گره دریافت تصویر و گره‌های پردازش استریو مورد استفاده قرار گرفت تا تمام مراحل تصحیح تصویر، تولید نقشه اختلاف مکانی و تخمین عمق بر اساس یک مجموعه پارامتر واحد انجام شوند.

\section{پارامترهای نهایی کالیبراسیون}
\label{app:stereo_calibration_final_params}

خلاصه پارامترهای اصلی حاصل از کالیبراسیون دوربین استریو در جدول \ref{tab:stereo_calibration_final_params} ارائه شده است. این پارامترها همان مقادیری هستند که در زنجیره پردازش تصویر برای تصحیح هندسی، تولید نقشه اختلاف مکانی و محاسبه عمق مورد استفاده قرار گرفتند.

\begin{table}[htbp]
	\centering
	\caption{خلاصه پارامترهای نهایی کالیبراسیون دوربین استریو}
	\label{tab:stereo_calibration_final_params}
	\renewcommand{\arraystretch}{1.45}
	\begin{tabular}{|c|c|}
		\hline
		\textbf{پارامتر} & \textbf{مقدار} \\
		\hline
		رزولوشن کالیبراسیون و پردازش تصویر & $1640 \times 1232$ پیکسل \\
		\hline
		تعداد گوشه‌های داخلی الگوی شطرنجی & $10 \times 7$ \\
		\hline
		اندازه هر خانه شطرنجی & $25~mm$ \\
		\hline
		فاصله پایه دوربین‌ها & $60.172~mm$ \\
		\hline
		فاصله پایه مورد استفاده در محاسبه عمق & $0.060172~m$ \\
		\hline
		فاصله کانونی تصویر تصحیح‌شده & $1238.075~px$ \\
		\hline
		مرکز تصویر تصحیح‌شده در راستای افقی & $785.625~px$ \\
		\hline
		مرکز تصویر تصحیح‌شده در راستای عمودی & $587.921~px$ \\
		\hline
		خطای بازفرافکنی استریو \lr{RMS} & $0.3774~px$ \\
		\hline
		ناحیه معتبر تصویر چپ & $1640 \times 1232$ پیکسل \\
		\hline
		ناحیه معتبر تصویر راست & $1640 \times 1231$ پیکسل \\
		\hline
	\end{tabular}
\end{table}

خطای بازفرافکنی استریو کمتر از یک پیکسل به دست آمد که برای کاربرد این پژوهش قابل قبول است. همچنین فاصله پایه محاسبه‌شده با فاصله فیزیکی ماژول دوربین استریو سازگار است. با توجه به رابطه محاسبه عمق در سامانه استریو، مقدار فاصله کانونی و فاصله پایه باید در تمام گره‌های پردازش تصویر از همین فایل کالیبراسیون خوانده شوند و از مقداردهی دستی یا استفاده از مقادیر قدیمی خودداری شود.

رابطه کلی تبدیل اختلاف مکانی به عمق به‌صورت زیر است:

\begin{equation}
	Z = \frac{fB}{d}
	\label{eq:appendix_stereo_depth_relation}
\end{equation}

در رابطه \ref{eq:appendix_stereo_depth_relation}، $Z$ عمق نقطه نسبت به دوربین، $f$ فاصله کانونی تصویر تصحیح‌شده، $B$ فاصله پایه دوربین‌ها و $d$ مقدار اختلاف مکانی است. بنابراین، دقت تخمین عمق به‌طور مستقیم به کیفیت کالیبراسیون، هم‌راستاسازی صحیح تصاویر و اعتبار مقدار اختلاف مکانی وابسته است.

\section{اعتبارسنجی تصحیح هندسی تصاویر}
\label{app:stereo_rectification_validation}

پس از محاسبه پارامترهای کالیبراسیون، صحت تصحیح هندسی تصاویر بررسی شد. برای این منظور، یک ابزار تشخیصی مستقل استفاده شد که جفت‌تصویر چپ و راست را دریافت کرده، نقشه‌های تصحیح هندسی را از فایل کالیبراسیون می‌سازد و سپس خطوط افقی متناظر را روی تصاویر تصحیح‌شده رسم می‌کند. اگر تصحیح هندسی به‌درستی انجام شده باشد، نقاط متناظر در دو تصویر باید روی خطوط افقی یکسان قرار گیرند.

علاوه بر بررسی خطوط متناظر، ترتیب دوربین‌های چپ و راست نیز ارزیابی شد. در این آزمون، دو حالت مختلف برای ترتیب دوربین‌ها بررسی می‌شود و برای هر حالت، هم تصویر دارای خطوط متناظر و هم نقشه اختلاف مکانی آزمایشی تولید می‌گردد. حالت صحیح، حالتی است که هم خطوط متناظر در دو تصویر هم‌راستا باشند و هم نقشه اختلاف مکانی دارای ساختار پیوسته و سازگار با عمق صحنه باشد. در صورت جابه‌جایی دوربین چپ و راست، ممکن است تصویر ظاهراً قابل مشاهده باشد، اما نقشه اختلاف مکانی ناپایدار، پراکنده و فاقد ساختار فیزیکی مناسب خواهد شد.

نمونه‌ای از خروجی آزمون خطوط متناظر و نقشه اختلاف مکانی در شکل \ref{fig:stereo_rectification_validation_outputs} نشان داده شده است. در تصاویر خطوط متناظر، خط‌های افقی روی تصویر چپ و راست باید از نقاط مشابه صحنه عبور کنند. در نقشه اختلاف مکانی نیز نواحی نزدیک‌تر باید به‌صورت پیوسته‌تر و با مقادیر اختلاف مکانی بیشتر ظاهر شوند، در حالی که نواحی دورتر اختلاف مکانی کمتری دارند.

\begin{figure}[htbp]
	\centering
	
	\begin{minipage}{0.48\linewidth}
		\centering
		\includegraphics[width=\linewidth]{epipolar_test1.png}
		\vspace{2mm}
		
		{\small آزمون خطوط متناظر، حالت اول}
	\end{minipage}
	\hfill
	\begin{minipage}{0.48\linewidth}
		\centering
		\includegraphics[width=\linewidth]{disparity_test1.png}
		\vspace{2mm}
		
		{\small نقشه اختلاف مکانی، حالت اول}
	\end{minipage}
	
	\vspace{4mm}
	
	\begin{minipage}{0.48\linewidth}
		\centering
		\includegraphics[width=\linewidth]{epipolar_test2.png}
		\vspace{2mm}
		
		{\small آزمون خطوط متناظر، حالت دوم}
	\end{minipage}
	\hfill
	\begin{minipage}{0.48\linewidth}
		\centering
		\includegraphics[width=\linewidth]{disparity_test2.png}
		\vspace{2mm}
		
		{\small نقشه اختلاف مکانی، حالت دوم}
	\end{minipage}
	
	\caption{اعتبارسنجی تصحیح هندسی تصاویر و بررسی ترتیب دوربین‌های چپ و راست}
	\label{fig:stereo_rectification_validation_outputs}
\end{figure}

این آزمون علاوه بر بررسی عددی خطای کالیبراسیون، یک ارزیابی کیفی مهم از صحت استفاده عملی از پارامترهای کالیبراسیون فراهم می‌کند. دلیل اهمیت این مرحله آن است که حتی با وجود خطای بازفرافکنی مناسب، عواملی مانند تغییر جهت تصویر، جابه‌جایی ترتیب دوربین‌ها، استفاده از رزولوشن متفاوت یا اعمال دوباره تصحیح هندسی می‌تواند باعث افت شدید کیفیت نقشه اختلاف مکانی شود. بنابراین، پس از کالیبراسیون، صحت هم‌راستاسازی تصاویر و سازگاری نقشه اختلاف مکانی به‌صورت جداگانه بررسی شد.

\section{نکات اجرایی و محدودیت‌ها}
\label{app:stereo_calibration_notes}

در اجرای کالیبراسیون و استفاده از پارامترهای آن، چند نکته اجرایی باید رعایت شود. نخست آنکه رزولوشن تصاویر مورد استفاده در پردازش باید با رزولوشن کالیبراسیون یکسان باشد. در این پژوهش، کالیبراسیون و پردازش تصویر هر دو با رزولوشن $1640 \times 1232$ انجام شدند. دوم آنکه جهت تصویر پس از کالیبراسیون نباید تغییر کند. هرگونه چرخش، وارون‌سازی یا برش نامنطبق با فایل کالیبراسیون می‌تواند باعث از بین رفتن هم‌راستاسازی تصاویر شود.

سوم آنکه تصحیح هندسی باید تنها در یک نقطه از زنجیره پردازش اعمال شود. در این پژوهش، تصحیح هندسی در گره دریافت تصویر انجام شده و گره تولید نقشه اختلاف مکانی تصاویر را دوباره تصحیح نمی‌کند. این تصمیم برای جلوگیری از اعمال دوباره تبدیل هندسی و حفظ سازگاری تصاویر اتخاذ شد.

چهارم آنکه کیفیت تصاویر ثبت‌شده برای کالیبراسیون اثر مستقیمی بر نتیجه دارد. تصاویر تار، تصاویر دارای بازتاب شدید، تصاویر با پوشش کم الگوی شطرنجی و تصاویر فاقد تنوع مکانی کافی باعث کاهش اعتبار پارامترهای کالیبراسیون می‌شوند. به همین دلیل، در هنگام ثبت داده‌ها تلاش شد الگوی شطرنجی در بخش‌های مختلف میدان دید و با زاویه‌ها و فاصله‌های متفاوت قرار گیرد.

در مجموع، فرایند کالیبراسیون دوربین استریو شامل دریافت تصاویر با رزولوشن ثابت، استخراج تصویر خاکستری از صفحه روشنایی، ثبت جفت‌تصویرهای متنوع از الگوی شطرنجی، آشکارسازی گوشه‌های داخلی، محاسبه پارامترهای درونی و بیرونی دوربین‌ها، تولید ماتریس‌های تصحیح هندسی، ذخیره خروجی در فایل کالیبراسیون و اعتبارسنجی نهایی با خطوط متناظر و نقشه اختلاف مکانی بود. نتایج حاصل نشان دادند که پارامترهای استخراج‌شده برای استفاده در زنجیره ادراک سامانه عملی مناسب هستند و می‌توانند مبنای تولید نقشه اختلاف مکانی و تخمین فاصله هدف و موانع قرار گیرند.